\newpage
\chapter{PROBLEM DEFINITION AND RESEARCH QUESTIONS}
\pagestyle{fancy}
\pagenumbering{arabic}

% ================================================
% 2.1 Introduction
% ================================================
\section{Introduction}

This chapter formalizes the research problem addressed by this thesis. Based on empirical analysis of the NYC Yellow Taxi dataset, we identify three critical gaps in existing streaming DQ frameworks, and derive five research questions that guide the rest of this work.

% ================================================
% 2.2 Empirical Analysis
% ================================================
\section{Empirical Analysis from NYC Taxi Dataset}

The empirical analysis reveals four fundamental patterns in the NYC Taxi dataset that motivate the CA-DQStream design.

\subsection{Dataset Overview}

Analyzing 3,076,903 records from July 2024 with baseline from January 2024 (2,964,624 records) reveals that 473,554 records (15.4\%) have at least one data quality violation. Critically, 84\% of violations are data completeness issues (NULL fields), not statistical outliers.

The 4-layer architecture is directly motivated by this finding: Layer 1 addresses completeness, Layer 2 addresses validity/consistency, Layer 3 addresses accuracy within context, and Layer 4 addresses multivariate anomalies not covered by rules.

\subsection{Layer 1 --- Schema Check Results}

The Schema Checker identifies completeness violations on required fields:

\begin{longtable}{|m{5cm}|R{2.5cm}|R{2cm}|m{3.5cm}|}
\hline
\textbf{Violation Type} & \textbf{Records} & \textbf{Rate} & \textbf{DQ Dimension} \\ \hline
NULL passenger\_count & 278,989 & 9.07\% & Completeness \\ \hline
NULL RatecodeID & 278,989 & 9.07\% & Completeness \\ \hline
NULL (simultaneous) & 278,989 & 9.07\% & Completeness \\ \hline
\end{longtable}

These records represent the same 278,989 records with both fields NULL. Since passenger count and ratecode are essential for downstream processing, these records are rejected at Layer 1.

\subsection{Layer 2 --- Hard Rules Results}

Hard Rules identify validity and consistency violations:

\begin{longtable}{|m{5cm}|R{2.5cm}|R{2cm}|m{3cm}|}
\hline
\textbf{Violation Type} & \textbf{Records} & \textbf{Rate} & \textbf{Threshold} \\ \hline
Negative fare\_amount & 58,650 & 1.91\% & fare\_amount $\geq$ 0 \\ \hline
Zero trip\_distance & 48,045 & 1.56\% & trip\_distance $>$ 0 \\ \hline
Extreme speed ($>$100 mph) & 2,064 & 0.07\% & speed $<$ 100 mph \\ \hline
Dropoff $\leq$ Pickup & 56 & $<$0.01\% & dropoff\_datetime $>$ pickup\_datetime \\ \hline
\end{longtable}

These rules are deterministic and require no context --- they apply universally.

\subsection{Layer 3 --- Context-Aware Rules Results}

Context-Aware Rules identify accuracy violations within specific context groups:

\begin{longtable}{|m{5cm}|R{2.5cm}|R{2cm}|m{3cm}|}
\hline
\textbf{Violation Type} & \textbf{Records} & \textbf{Rate} & \textbf{Method} \\ \hline
fare\_outlier ($|x-\mu|/\sigma > 2$) & 80,209 & 2.61\% & Context z-score \\ \hline
distance\_outlier ($|x-\mu|/\sigma > 2$) & 75,255 & 2.45\% & Context z-score \\ \hline
total\_outlier ($|x-\mu|/\sigma > 2$) & 98,300 & 3.19\% & Context z-score \\ \hline
\end{longtable}

Context groups are defined by the combination of time slot (night, morning, afternoon, evening) and ratecode type (standard, special). This separation is critical: special ratecode trips (RatecodeID $\neq$ 1) have average fares 2--3$\times$ higher than standard trips in the same time slot, and variance up to 20$\times$ larger.

\begin{longtable}{|m{3.5cm}|m{2.5cm}|m{2.5cm}|m{2.5cm}|m{2cm}|}
\hline
\textbf{Context Group} & \textbf{Avg Fare (\$)} & \textbf{Std Dev} & \textbf{Threshold} & \textbf{Count} \\ \hline
night standard & 15.64 & 13.34 & 42.32 & 191,819 \\ \hline
night special & 36.25 & 32.98 & 102.21 & 38,736 \\ \hline
morning standard & 15.26 & 12.55 & 40.36 & 588,251 \\ \hline
morning special & 40.26 & 44.20 & 128.66 & 72,396 \\ \hline
afternoon standard & 15.31 & 12.03 & 39.37 & 1,002,632 \\ \hline
afternoon special & 49.12 & 36.98 & 123.08 & 100,545 \\ \hline
evening standard & 15.44 & 12.37 & 40.18 & 880,648 \\ \hline
evening special & 41.30 & 37.16 & 115.62 & 89,597 \\ \hline
\end{longtable}

The fixed threshold approach (e.g., p95 = $\$61.8$ from January 2024) would over-flag nearly 100\% of special trips. The context-aware threshold of $\$115$--$\$129$ for special ratecode correctly identifies only the true outliers.

\begin{figure}[H]
\centering
\includegraphics[width=0.85\textwidth]{figs/F3_normal_vs_special.png}
\caption{Comparison of fare and distance distributions between Standard (RatecodeID$=1$) and Special (RatecodeID$\neq$1) ratecode groups on January 2024 data. Special trips have significantly higher median fares and greater variance, justifying the context-aware approach.}
\label{fig:f3_normal_vs_special}
\end{figure}

\begin{figure}[H]
\centering
\includegraphics[width=0.85\textwidth]{figs/T7_normal_vs_special_copy.png}
\caption{Standard vs Special: volume and average fare by time bin. Special trips are consistently 2--3$\times$ higher in average fare across all time slots, demonstrating the need for context-specific thresholds.}
\label{fig:t7_normal_vs_special}
\end{figure}

% ================================================
% 2.3 Layer 4 Results
% ================================================
\section{Layer 4 --- ML Layer Scope}

Rules (Layers 1--3) collectively cover approximately 97\% of all violations. Layer 4 focuses on the remaining 3\% of cases where single-feature rules are insufficient:

\begin{itemize}
    \item \textbf{Multivariate combination anomalies:} fare, distance, and duration jointly anomalous, but each feature individually within threshold.
    \item \textbf{Subtle ratio anomalies:} fare-to-distance ratio exceeds normal range (e.g., $\$80$ fare for 0.5 miles, where both fare$>$0 and distance$>$0 pass individually).
    \item \textbf{Zone-level aggregate drift:} zone average fare shifts without individual records being statistical outliers.
    \item \textbf{Cross-feature inconsistencies:} payment type and tip amount contradiction at boundary cases.
\end{itemize}

\begin{figure}[H]
\centering
\subfloat[Distance vs Fare by VendorID (Jan 2024)]{
\label{fig:fig5}
\includegraphics[width=0.48\textwidth]{figs/5_Distance_vs_Fare.png}
}
\subfloat[Distance vs Duration by VendorID (Jan 2024)]{
\label{fig:fig6}
\includegraphics[width=0.48\textwidth]{figs/6_Distance_vs_Duration.png}
}
\caption{Scatter plots of fare vs. distance and distance vs. duration by VendorID. The linear relationship (slope $\approx$ \$2.5/mile) provides a reference for detecting multivariate anomalies.}
\end{figure}

% ================================================
% 2.4 Research Gaps
% ================================================
\section{Research Gaps}

Based on the empirical analysis, three research gaps are identified:

\subsection{Gap 1: No Formal Context Model}

No existing framework defines context formally. Stream DaQ (2025) supports dynamic adaptation but lacks mathematical definitions. The empirical evidence is clear: fixed threshold $\$61.8$ (p95 from training) over-flags 99\% of special ratecode trips because special ratecode has max fare $\$2{,}500$--$\$5{,}000$ but p95 only $\$82$--$\$88$. Without a context model, FPR is inevitably high.

\subsection{Gap 2: No ML Integration for Streaming DQ}

No open-source framework integrates ML-based anomaly detection for streaming DQ. Uber and Grab use Isolation Forest and Random Cut Forest in production, but implementations are proprietary. The multivariate anomalies that rules cannot express remain undetected.

\subsection{Gap 3: No Self-Monitoring of the DQ System}

No framework monitors its own output to detect when baselines have become stale. NULL rates are observed to increase from 2024 to 2026 (evidenced by monthly violation trends). Without meta-level monitoring, gradual concept drift goes undetected until the DQ model produces unreliable results.

\begin{figure}[H]
\centering
\subfloat[Monthly Record Volume (Jan 2024 -- Feb 2026)]{
\label{fig:fig1}
\includegraphics[width=0.48\textwidth]{figs/1_Monthly_Record_Volume.png}
}
\subfloat[Violation Breakdown by Rule (Jan 2024)]{
\label{fig:fig2}
\includegraphics[width=0.48\textwidth]{figs/2_Violation_Breakdown.png}
}
\caption{(Left) Monthly record volume showing seasonal variations of 2--3$\times$ between peak and off-peak months. January 2024 (training) and July 2024 (test drift) are highlighted. (Right) Violation breakdown by rule type, with NULL passenger and ratecode violations dominating at 278,989 records each.}
\end{figure}

\begin{figure}[H]
\centering
\subfloat[Violation Rate by VendorID across months]{
\label{fig:fig3}
\includegraphics[width=0.65\textwidth]{figs/3_Violation_Rate_by_VendorID.png}
}
\caption{Violation rates vary significantly by VendorID. Vendor 6 (Unknown) consistently shows higher violation rates, and Vendor 2 shows distinct seasonal patterns. This vendor-specific variation is captured by context-aware rules but would be missed by a single global threshold.}
\end{figure}

\begin{figure}[H]
\centering
\subfloat[Violation counts over time (Jan 2024 -- Feb 2026)]{
\label{fig:t5}
\includegraphics[width=0.85\textwidth]{figs/T5_violations_by_month.png}
}
\caption{Stacked violation counts over 26 months. The stacked area shows total violations growing with data volume, while the bottom panel reveals the violation rate trends. NULL passenger violations show a gradual increase from 2024 to 2026 --- evidence of gradual concept drift requiring ADWIN-based monitoring.}
\end{figure}

% ================================================
% 2.5 Research Questions
% ================================================
\section{Research Questions}

This thesis addresses five research questions organized by theme:

\subsection{RQ1: Formal Context Model}
\textit{``How can context be formally defined for streaming DQ, and does context separation justify the overhead?''}

\begin{itemize}
    \item \textbf{RQ1a:} Can the Context Separation Score (CSS) exceed 2.0 for all time slots? CSS = $T_{special} / T_{standard}$ where $T = \mu + 2\sigma$. If CSS $>$ 2.0, special vs. standard separation is statistically justified.
    \item \textbf{RQ1b:} How does the hierarchical fallback mechanism (2D $\to$ 1D $\to$ global) perform in practice?
    \item \textbf{RQ1c:} Is $n_{min} = 100$ samples sufficient for reliable baseline computation?
\end{itemize}

\subsection{RQ2: ML Anomaly Detection}
\textit{``Can ML methods detect multivariate anomalies that rules miss, and which method performs best?''}

\begin{itemize}
    \item \textbf{RQ2a:} What is the Precision/Recall of ML-only anomalies (records not flagged by Layers 1--3)?
    \item \textbf{RQ2b:} Does online retraining via partial\_fit() maintain ML performance over time?
    \item \textbf{RQ2c:} Can zone-level aggregate z-score detect drift that context-aware rules miss?
    \item \textbf{RQ2d:} Which of Isolation Forest, LOF, or One-Class SVM performs best on Precision/Recall/F1/FPR?
    \item \textbf{RQ2e:} Does ensemble voting improve FPR over the best individual method?
\end{itemize}

\subsection{RQ3: Context-Aware vs. Fixed Rules}
\textit{``Does the context-aware approach significantly reduce false positive rate?''}

\begin{itemize}
    \item \textbf{RQ3a:} What is the FPR of context-aware rules on the NYC Taxi dataset? Target: $<$5\%.
    \item \textbf{RQ3b:} How does CA-DQStream compare to Great Expectations and Stream DaQ on the same dataset?
    \item \textbf{RQ3c:} How does detection latency compare between layers?
\end{itemize}

\subsection{RQ4: Drift Detection Effectiveness}
\textit{``Does ADWIN-based drift detection on quality meta-streams detect degradation promptly?''}

\begin{itemize}
    \item \textbf{RQ4a:} What is the Mean Time To Detect (MTTD) for gradual NULL rate increases?
    \item \textbf{RQ4b:} What is the false positive rate of the drift detector?
    \item \textbf{RQ4c:} Does self-adaptation (baseline recalculation) improve FPR after drift?
\end{itemize}

\subsection{RQ5: Timeliness and Uniqueness}
\textit{``Does meta-level monitoring detect timeliness and uniqueness issues?''}

\begin{itemize}
    \item \textbf{RQ5a:} Does consumer lag correlate with violation rate?
    \item \textbf{RQ5b:} How effective is Flink State-based deduplication for streaming?
    \item \textbf{RQ5c:} Do dedup rate spikes correlate with upstream retry behavior?
\end{itemize}

% ================================================
% 2.6 Research Hypotheses
% ================================================
\section{Research Hypotheses}

Based on the empirical analysis, we state the following hypotheses:

\begin{enumerate}
    \item \textbf{H1:} The Context Separation Score exceeds 2.0 for all time slots, justifying the 2D context model.
    \item \textbf{H2:} Context-aware thresholds achieve FPR below 5\% compared to 25--50\% for fixed thresholds.
    \item \textbf{H3:} ML layer detects at least 1\% of total violations not caught by Layers 1--3.
    \item \textbf{H4:} ADWIN detects gradual NULL rate drift within one monitoring window.
    \item \textbf{H5:} Self-adaptation (baseline recalculation) reduces FPR to pre-drift levels within two windows.
\end{enumerate}

% ================================================
% 2.7 Chapter Summary
% ================================================
\section{Chapter Summary}

This chapter formalized the research problem based on empirical analysis of the NYC Taxi dataset. The empirical evidence motivates the 4-layer architecture: Layer 1 addresses completeness (84\% of violations), Layer 2 addresses validity/consistency, Layer 3 addresses accuracy with context-aware thresholds, and Layer 4 handles multivariate anomalies. Three research gaps were identified: lack of formal context model, absence of ML integration, and no self-monitoring capability. Five research questions and five hypotheses guide the evaluation in Chapter 4. The next chapter details the CA-DQStream framework design.
